\chapter{Introduction}
\label{Chapter1}
\lhead{Chapter 1. \emph{Introduction}}

\section{Background and Motivation}

The last decade has witnessed an extraordinary acceleration in the capabilities of large language models (LLMs), enabling them to tackle increasingly specialised reasoning tasks across diverse domains, including medicine. Modern LLMs exhibit a high degree of linguistic fluency, contextual awareness, and domain-specific recall, which collectively position them as promising computational tools in clinical decision support. Their ability to synthesise medical guidelines, recall diagnostic criteria, and articulate differential diagnoses has generated significant interest in their potential integration into real-world healthcare workflows.

Nevertheless, it is widely recognised that clinical reasoning is fundamentally different from answering a static medical question. Real clinical encounters unfold as \textit{interactive}, \textit{sequential}, and \textit{uncertain} processes wherein physicians iteratively inquire, evaluate, hypothesise, refine, and re-evaluate based on newly acquired evidence. A typical clinical workflow consists of several cognitively demanding stages such as targeted history-taking, focused physical examination, ordering of laboratory and imaging tests, synthesis of multimodal evidence, and probabilistic reasoning under incomplete information. This iterative, exploratory process stands in sharp contrast to the rigid, one-shot question formats dominating traditional benchmarks such as MedQA or MedMCQA.

\section{Limitations of Static Benchmarks}

Widely utilised benchmarks such as MedQA, USMLE step examinations, and PubMedQA evaluate a model's ability to recall facts and apply clinical guidelines to fully structured vignettes. In these static tests, all necessary patient history, vital signs, and investigation results are provided upfront in a single prompt. While these benchmarks are useful for assessing baseline medical knowledge, they fail to evaluate the \textit{procedural and interactive} nature of clinical decision-making. They assume that the model already possesses all relevant patient information, thereby eliminating the need to ask clarifying questions, identify confounding factors, or selectively acquire evidence. As a result, such evaluations do not capture an LLM's capacity for hypothesis formation, adaptive reasoning, or multi-turn dialogue management---capabilities that are essential to real clinical practice.

\section{Evolution of the Research}

To bridge this evaluation gap, this thesis adopts and substantially extends the multi-agent LLM simulation paradigm.

\textbf{Initial Framework:} The foundational work laid the groundwork by implementing a functional multi-agent clinical simulator governed by an ad-hoc \texttt{while}-loop. Several open-source LLMs (Mistral-7B, Microsoft BioGPT, JSL-Med, and Apollo) were benchmarked and baseline metrics for diagnostic accuracy and bias susceptibility were established. However, this initial approach revealed severe architectural limitations—notably non-deterministic state management and a vulnerability termed \textit{Lexical Leakage}—arising from the deployment of homogeneous model instances across all agent roles.

\textbf{Enhanced Framework -- AgentClinic v2.1:} To achieve greater rigour and ecological validity, the system was entirely re-architected into \textbf{AgentClinic v2.1}. This evolution transitioned the orchestration layer to a LangGraph Directed Cyclic Graph (DCG) state machine, introduced 4-bit NF4 double quantisation to support heterogeneous multi-engine configurations on a single GPU, and formalised mathematical process-oriented metrics derived from a novel latent metacognitive reflection node.

\section{Problem Statement}

Despite notable progress in healthcare-focused LLM research, several persistent gaps remain:
\begin{itemize}
    \item \textbf{Absence of flexible, interactive frameworks} for evaluating open-source LLMs in realistic diagnostic workflows without relying on proprietary, high-bandwidth APIs (e.g., GPT-4).
    \item \textbf{Lexical Leakage} in existing open-source multi-agent systems, where the use of a single set of model weights for both the Doctor and Patient agents allows the Doctor to mathematically infer patient ground-truth representations through shared embedding spaces.
    \item \textbf{Insufficient process evaluation}, which focuses only on binary diagnostic accuracy rather than the cognitive trajectory, stability, and rationality of the LLM's reasoning process.
    \item \textbf{Lack of controlled bias modelling tools} to systematically evaluate how open-source LLMs mimic human cognitive errors, such as confirmation and recency biases.
\end{itemize}

\section{Contributions}

The primary contributions of this work include:
\begin{enumerate}
    \item A comprehensive analysis of baseline open-source LLM performance in interactive clinical settings using an ad-hoc while-loop simulation architecture, revealing the phenomenon of Lexical Leakage.
    \item The theoretical identification and mitigation of \textit{Lexical Leakage} via a heterogeneous multi-engine architecture utilising BitsAndBytes NF4 quantisation.
    \item A custom-built, LangGraph-orchestrated state machine featuring a \textit{latent metacognitive reflection node} to capture reasoning trajectories without patient-facing contamination.
    \item The mathematical formulation of three process-oriented clinical metrics: Diagnostic Stability (DS), Test Rationality (TR), and Information Efficiency (IE).
    \item A rigorous experimental suite demonstrating that structured metacognition significantly improves open-source baseline accuracy, alongside theoretically grounded projections for domain-specialised models and cognitive bias vulnerabilities.
\end{enumerate}

\section{Thesis Organisation}

This thesis is structured as follows. Chapter~2 reviews the relevant literature on medical AI benchmarking, clinical simulation, and cognitive bias. Chapter~3 describes the dataset and the simulated OSCE examination structure. Chapter~4 details the foundational framework and baseline results. Chapter~5 presents the enhanced methodology and full architecture of AgentClinic~v2.1. Chapter~6 reports the experimental results and projections. Chapter~7 provides a comprehensive discussion of the findings, and Chapter~8 concludes the thesis with directions for future research.
